\chapter{Introduction}\label{chap:introduction}
% \noindent You must use \textbackslash{}noindent at the beginning of the first paragraph in sections and subsections.

%

\begin{reviewed}

	This chapter is meant to contextualize the work explored and tackled throughout the dissertation by precisely defining the problem statement, the intended objectives, and the initial motivation for picking this area.

	\section{Motivation}

	Up until now, every technological achievement merely represented an extension of our capabilities. Now, with \gls{ai}, for the first time there is the prospect of a technology that would have its own goals and be capable of substituting the inventive process itself, rendering human capabilities obsolete.

	The motivation of this work stems from the desire to keep the future human and human aligned.

	\section{Problem}

	% Adiciona citações
	
	The dangers of unaligned \gls{ai} frequently start with Bostrom’s “Paperclip Maximizer” which demonstrates how an \gls{ai} system with seemingly benign goals could still potentially lead to catastrophic outcomes if not properly aligned with human values. This thought experiment is rooted in the orthogonality thesis which states that there is no inherent link between a system's goals and its ability to reach them.
	Building on these concerns, Yampolskiy argues that it's definitionally impossible for humans to understand or control a super-intelligent agent\cite{yampolskiy_controllability_2022}. Such arguments have motivated calls for caution in the development of advanced \gls{ai} systems, including public statements advocating for the suspension or regulation of certain forms of \gls{ai} research until adequate safety guarantees can be established\cite{FLI_SuperintelligenceStatement2025}.

	Today we have Multi-modal and Generalist \glspl{llm} capable of doing tasks in a large range of domains that exhibit early signs of instrumental convergence behavior\cite{greenblatt_alignment_faking}.

	Furthermore, current societal structure mandates, through competition and rivalry, that there will be large amounts of money invested in trying to reach such capable systems, and there are rising concerns that \gls{ai} companies have too great financial incentives to avoid effective self-regulating oversight\cite{hemphill_advanced_2025}. Witch is corroborated by several \gls{ai} safety reports that reveal how lacking the approach of current \gls{ai} companies is towards \gls{ai} safety\cite{FLIAISafetyIndex2025}\cite{AISafetyReport2025}.

	% Maybe explain how there is good work being done in grey literature https://www.alignment.org/blog/competing-with-sampling/ and how it's hard to parse for a normal person and how there's a bridge lacking between a normal person and real advancements in \gls{ai} safety research.

	% Knowing the gab between safety progress and \gls{ai} capabilities is widening I intend with my research to build on work such as: \gls{ai} Control: Improving Safety Despite Intentional Subversion. My goal is to mitigate or demonstrate the risks of unsafe systems, in the hope that sufficiently robust misalignment demonstrations may instill enough public discourse to warrant and enable further systemic action.

\end{reviewed}

\begin{toreview}

\section{Objectives}
The world of \gls{ai} safety is in and of itself vast and unfocused, large amounts of research on very different areas exist. To steer the research objectives, a hypothesis is proposed:
\begin{itemize}
	\item \textbf{Hypothesis}: If the current trajectory of \gls{ai} development is deemed to be unsafe. Then, in the present environment where large incentives to invest money into this technology exist, \gls{ai} safety research alone can't be the silver bullet. Thus, to truly guarantee safe and aligned \gls{ai} systems for humans in the future, structural changes are also required.
\end{itemize}
Therefore, this research will be conducted in a \gls{slr} that aims at finding out if the current pursuit of \gls{agi} has natural dangers, how dangerous, and whether they are relevant in general or only narrow domains.

Thus, this literature review shall focus on work that has been carried out in analyzing the intent and inner objectives of current generalist and narrow models. For that purpose, the following research questions will be used to carry out the research on the \textbf{SLR}:

\begin{enumerate}
	\item \gls{rq}1. What is the main focus of current \gls{ai} safety research?
	\item \gls{rq}2. Which are the current methods used to look into the alignment of current generalist models?
	\item \gls{rq}3. What conclusions regarding the safety of current models and safety of development and deployment processes can we assert from current research?
	\item \gls{rq}4. Is there a measurable gap between the research attention given to an AI safety area and its assessed level of danger?
\end{enumerate}

After conclusions are drawn from these research questions a path for the most relevant \gls{ai} safety research field to explore can be laid out.

\end{toreview}

\section{Dissertation Structure}

\newpage
